\documentclass[zad,zawodnik,utf8,pl]{sinol}
\title{TODO}
\id{__ID__}
\signature{TODO}
\konkurs{TODO}
\day{TODO}
\date{TODO}
\RAM{512}
\iomode{interactive}
\pagestyle{fancy}
\setlength{\parindent}{0pt}  % no paragraph indent
\setlength{\parskip}{0.6em}  % visible gap between paragraphs instead

\begin{document}
\begin{tasktext}

\vspace{-0.5cm}  % smaller gap above the first paragraph
\noindent {}%
% TEMPLATE: write the story and the task here.
TODO

\section{Komunikacja}
% Describe the protocol exactly: who writes first, the format of every
% message, the query limit, and that every line must be flushed. Also what is in
% dlazaw.zip (the sample interactor, tests) and how to run it.
TODO

Po wypisaniu każdego wiersza Twój program musi opróżnić bufor wyjścia,
np.\ \texttt{cout << endl}, \texttt{fflush(stdout)} lub \texttt{cout.flush()}.

\section{Przykład}
% Interactive tasks have no .out files - show an example interaction by hand.
\begin{center}
\begin{tabular}{|l|l|}
\hline
\textbf{Program zawodnika} & \textbf{Biblioteka} \\ \hline
 & \texttt{TODO} \\ \hline
\texttt{TODO} & \texttt{TODO} \\ \hline
\end{tabular}
\end{center}

\section{Podzadania}
% Keep in sync with `scores` and `time_limit(s)` in config.yml (checked by `tools/contest.py lint`).
\begin{center}
\begin{tabular}{|l|l|l|l|}
\hline
\textbf{\sinolSNameSubtask} & \textbf{\sinolSNamePoints} &
\textbf{\sinolSNameConditions} & \textbf{\sinolSNameTimeLimit} \\ \hline
1 & $30$ & $n \le 20$ & $1$ s \\ \hline
2 & $70$ & brak dodatkowych założeń & $1$ s \\ \hline
\end{tabular}
\end{center}

\end{tasktext}
\end{document}
